\documentclass[11pt,a4paper]{article}
\usepackage[margin=1in]{geometry}
\usepackage{amsmath,amssymb,amsthm}
\usepackage{algorithm}
\usepackage{algpseudocode}
\usepackage{booktabs}
\usepackage{hyperref}

\theoremstyle{definition}
\newtheorem{definition}{Definition}[section]
\newtheorem{theorem}{Theorem}[section]
\newtheorem{lemma}[theorem]{Lemma}
\newtheorem{remark}{Remark}[section]

\title{\textbf{Theoretical Foundations, Checkerboard Artifact Dynamics, and Spatial Rearrangement of Transposed Convolution and Pixel Shuffle}}
\author{Team Scaibu \\ \small Email: \texttt{jaydeep.9664920749@gmail.com} \\ \small Architecture and Core Engine Team}
\date{October 2026}

\begin{document}
\maketitle

\begin{abstract}
This document provides the formal mathematical specification, dual-gradient transpose mechanics, and sub-pixel periodic rearrangement theory for Transposed Convolution and Pixel Shuffle (Sub-Pixel Convolution). Generative architectures, autoencoders, and super-resolution models require spatial magnification of low-resolution feature maps. Transposed convolution implements the gradient adjoint of a strided convolution, frequently introducing high-frequency checkerboard artifacts due to uneven overlap. Pixel shuffle circumvents artifacts by computing $r^2 C$ lower-resolution channels and periodically interleaving them into an $r\times$ magnified grid. We derive the output spatial equations, prove checkerboard-free reconstruction guarantees, analyze computational complexities, trace a worked numerical example, and document memory alignment practices.
\end{abstract}

\section{Notation and Mathematical Preliminaries}

Let $X \in \mathbb{R}^{C_{\text{in}} \times H_{\text{in}} \times W_{\text{in}}}$ denote an input low-resolution activation tensor.

\begin{itemize}
    \item $C_{\text{in}}, C_{\text{out}} \in \mathbb{N}_{\ge 1}$: Input and output channel counts.
    \item $H_{\text{in}}, W_{\text{in}} \in \mathbb{N}_{\ge 1}$: Input spatial grid resolution.
    \item $r \in \mathbb{N}_{\ge 1}$: Spatial magnification upscaling ratio ($r = 2$).
    \item $W \in \mathbb{R}^{C_{\text{in}} \times C_{\text{out}} \times K_h \times K_w}$: Transposed convolution filter tensor.
    \item $s \in \mathbb{N}_{\ge 1}, p \in \mathbb{N}_{\ge 0}$: Stride and zero-padding factors for transposed convolution.
    \item $Y \in \mathbb{R}^{C_{\text{out}} \times H_{\text{out}} \times W_{\text{out}}}$: Upsampled high-resolution output tensor.
\end{itemize}

\begin{table}[h]
\centering
\caption{Summary of Mathematical Symbols for Spatial Upsampling}
\begin{tabular}{lll}
\toprule
\textbf{Symbol} & \textbf{Type / Domain} & \textbf{Description} \\
\midrule
$X$ & $\mathbb{R}^{C_{\text{in}} \times H_{\text{in}} \times W_{\text{in}}}$ & Input low-resolution feature tensor \\
$r$ & $\mathbb{N}_{\ge 1}$ & Pixel shuffle upscaling factor \\
$W$ & $\mathbb{R}^{C_{\text{in}} \times C_{\text{out}} \times K_h \times K_w}$ & Transposed convolution kernel \\
$H_{\text{out}}, W_{\text{out}}$ & $\mathbb{N}_{\ge 1}$ & High-resolution spatial dimensions \\
$Y$ & $\mathbb{R}^{C_{\text{out}} \times H_{\text{out}} \times W_{\text{out}}}$ & Upsampled output tensor \\
\bottomrule
\end{tabular}
\end{table}

\section{Problem Definition and Core Concepts}

\subsection{Intuitive Meaning}
While standard convolution reduces spatial dimension through striding, upsampling expands low-resolution feature maps into high-resolution images. Transposed convolution scatters each input pixel into a $K_h \times K_w$ output patch scaled by filter weights. Pixel Shuffle computes extra channels at low resolution and rearranges them into higher spatial resolution, eliminating uneven stride overlap artifacts.

\subsection{Formal Mathematical Formulations}

\subsubsection{Transposed Convolution (Fractionally Strided Convolution)}
The output resolution is:
\begin{align}
H_{\text{out}} &= (H_{\text{in}} - 1) \cdot s - 2p + K_h \\
W_{\text{out}} &= (W_{\text{in}} - 1) \cdot s - 2p + K_w
\end{align}
For each input coordinate $(h_{\text{in}}, w_{\text{in}})$:
\begin{equation}
Y(c_{\text{out}}, h_{\text{in}} s - p + i, w_{\text{in}} s - p + j) += \sum_{c_{\text{in}}=0}^{C_{\text{in}}-1} X(c_{\text{in}}, h_{\text{in}}, w_{\text{in}}) \cdot W(c_{\text{in}}, c_{\text{out}}, i, j)
\end{equation}

\subsubsection{Pixel Shuffle (Periodic Sub-Pixel Interleaving)}
With $C_{\text{in}} = r^2 C_{\text{out}}$, for offsets $v, u \in \{0, \dots, r-1\}$:
\begin{equation}
Y(c_{\text{out}}, h_{\text{in}} r + v, w_{\text{in}} r + u) = X(c_{\text{out}} r^2 + v r + u, h_{\text{in}}, w_{\text{in}})
\end{equation}

\section{Algorithm Specification}

\begin{algorithm}[H]
\caption{Transposed Convolution and Pixel Shuffle Upsampling}
\label{alg:transposed_conv}
\begin{algorithmic}[1]
\Require Input $X \in \mathbb{R}^{C_{\text{in}} \times H_{\text{in}} \times W_{\text{in}}}$, mode $M \in \{\text{transposed\_conv}, \text{pixel\_shuffle}\}$, weights $W$, stride $s$, pad $p$, upscale $r$.
\Ensure Upsampled tensor $Y \in \mathbb{R}^{C_{\text{out}} \times H_{\text{out}} \times W_{\text{out}}}$, dimensions $H_{\text{out}}, W_{\text{out}}$.
\If{$M = \text{pixel\_shuffle}$}
    \State $C_{\text{out}} \gets C_{\text{in}} / r^2$, $H_{\text{out}} \gets H_{\text{in}} \cdot r$, $W_{\text{out}} \gets W_{\text{in}} \cdot r$
    \State Allocate $Y \in \mathbb{R}^{C_{\text{out}} \times H_{\text{out}} \times W_{\text{out}}}$
    \For{$c_{\text{out}} \gets 0 \textbf{ to } C_{\text{out}} - 1$}
        \For{$h \gets 0 \textbf{ to } H_{\text{in}} - 1$}
            \For{$w \gets 0 \textbf{ to } W_{\text{in}} - 1$}
                \For{$v \gets 0 \textbf{ to } r - 1$}
                    \For{$u \gets 0 \textbf{ to } r - 1$}
                        \State $c_{\text{in}} \gets c_{\text{out}} r^2 + v r + u$
                        \State $Y[c_{\text{out}}][h \cdot r + v][w \cdot r + u] \gets X[c_{\text{in}}][h][w]$
                    \EndFor
                \EndFor
            \EndFor
        \EndFor
    \EndFor
\Else
    \State $C_{\text{out}} \gets \text{dim}_1(W)$, $H_{\text{out}} \gets (H_{\text{in}} - 1)s - 2p + K_h$, $W_{\text{out}} \gets (W_{\text{in}} - 1)s - 2p + K_w$
    \State Allocate $Y \gets \mathbf{0}^{C_{\text{out}} \times H_{\text{out}} \times W_{\text{out}}}$
    \For{$c_{\text{in}} \gets 0 \textbf{ to } C_{\text{in}} - 1$}
        \For{$c_{\text{out}} \gets 0 \textbf{ to } C_{\text{out}} - 1$}
            \For{$h \gets 0 \textbf{ to } H_{\text{in}} - 1$}
                \For{$w \gets 0 \textbf{ to } W_{\text{in}} - 1$}
                    \For{$i \gets 0 \textbf{ to } K_h - 1$}
                        \For{$j \gets 0 \textbf{ to } K_w - 1$}
                            \State $Y[c_{\text{out}}][h s - p + i][w s - p + j] \mathrel{+}= X[c_{\text{in}}][h][w] \cdot W[c_{\text{in}}][c_{\text{out}}][i][j]$
                        \EndFor
                    \EndFor
                \EndFor
            \EndFor
        \EndFor
    \EndFor
\EndIf
\State \Return $\{Y, H_{\text{out}}, W_{\text{out}}\}$
\end{algorithmic}
\end{algorithm}

\section{Theoretical Guarantees and Artifact Elimination}

\begin{theorem}[Checkerboard Elimination Guarantee of Pixel Shuffle]
In 2D transposed convolution, if kernel size $K$ is not divisible by stride $s$ ($K \pmod s \ne 0$), the number of overlapping filter contributions across spatial coordinates varies periodically, inducing 2D checkerboard artifacts. In Pixel Shuffle, every output pixel receives an exact, uniform bijective mapping from the feature channels:
\begin{equation}
|\{(c_{\text{in}}, h_{\text{in}}, w_{\text{in}}) \to (c_{\text{out}}, h_{\text{out}}, w_{\text{out}})\}| = 1 \quad \forall (c_{\text{out}}, h_{\text{out}}, w_{\text{out}})
\end{equation}
guaranteeing complete mathematical elimination of uneven overlap checkerboard patterns.
\end{theorem}

\begin{proof}
Because the mapping $(c_{\text{out}}, h_{\text{in}}, w_{\text{in}}, v, u) \mapsto (c_{\text{out}}, h_{\text{in}} r + v, w_{\text{in}} r + u)$ is a bijection between $\mathbb{Z}_{C_{\text{out}}} \times \mathbb{Z}_{H_{\text{in}}} \times \mathbb{Z}_{W_{\text{in}}} \times \mathbb{Z}_r \times \mathbb{Z}_r$ and $\mathbb{Z}_{C_{\text{out}}} \times \mathbb{Z}_{H_{\text{in}} r} \times \mathbb{Z}_{W_{\text{in}} r}$, each output grid position receives exactly one feature value with zero spatial overlap.
\end{proof}

\subsection{Limitations}
\begin{itemize}
    \item \textbf{Channel Dimensionality Requirement}: Pixel Shuffle strictly requires input channels $C_{\text{in}}$ to be an exact integer multiple of $r^2$.
\end{itemize}

\section{Complexity Analysis}

\subsection{Time Complexity}
\begin{itemize}
    \item \textbf{Pixel Shuffle}: Zero floating-point arithmetic multiplications; executes pure memory address permuting in $\Theta(C_{\text{in}} \cdot H_{\text{in}} \cdot W_{\text{in}})$ memory copies.
    \item \textbf{Transposed Convolution}: $\Theta(C_{\text{in}} \cdot C_{\text{out}} \cdot H_{\text{in}} \cdot W_{\text{in}} \cdot K_h \cdot K_w)$ FLOPs.
\end{itemize}

\subsection{Space Complexity}
\begin{itemize}
    \item \textbf{Storage}: $\mathcal{O}(C_{\text{out}} \cdot H_{\text{out}} \cdot W_{\text{out}})$ memory for the upsampled output.
\end{itemize}

\section{Exactness and Error Analysis}

Pixel shuffle is an exact bijective tensor permutation without numerical precision loss.

\section{Worked Numerical Example}

Let $C_{\text{in}} = 4, H_{\text{in}} = 1, W_{\text{in}} = 1, r = 2 \implies C_{\text{out}} = 1, H_{\text{out}} = 2, W_{\text{out}} = 2$.
\begin{equation}
X = [[[10.0]], [[20.0]], [[30.0]], [[40.0]]] \in \mathbb{R}^{4 \times 1 \times 1}
\end{equation}

\begin{enumerate}
    \item Pixel Shuffle mapping ($r = 2$):
    \begin{align}
    (v=0, u=0): c_{\text{in}} = 0 \cdot 4 + 0 \cdot 2 + 0 = 0 &\implies Y(0, 0, 0) = X(0, 0, 0) = 10.0 \\
    (v=0, u=1): c_{\text{in}} = 0 \cdot 4 + 0 \cdot 2 + 1 = 1 &\implies Y(0, 0, 1) = X(1, 0, 0) = 20.0 \\
    (v=1, u=0): c_{\text{in}} = 0 \cdot 4 + 1 \cdot 2 + 0 = 2 &\implies Y(0, 1, 0) = X(2, 0, 0) = 30.0 \\
    (v=1, u=1): c_{\text{in}} = 0 \cdot 4 + 1 \cdot 2 + 1 = 3 &\implies Y(0, 1, 1) = X(3, 0, 0) = 40.0
    \end{align}
    Output: $Y = [[[10.0, 20.0], [30.0, 40.0]]] \in \mathbb{R}^{1 \times 2 \times 2}$.
\end{enumerate}

\section{Numerical Considerations and Edge Cases}

\begin{itemize}
    \item \textbf{Memory Contiguity}: Output tensor layout corresponds to standard PyTorch/TensorFlow NCHW and NHWC memory packing without requiring intermediate array buffering.
\end{itemize}

\begin{thebibliography}{9}
\bibitem{shi2016real} W.~Shi et al., ``Real-Time Single Image and Video Super-Resolution Using an Efficient Sub-Pixel Convolutional Neural Network,'' in \emph{IEEE Conference on Computer Vision and Pattern Recognition (CVPR)}, pp.~1874--1883, 2016.
\bibitem{dumoulin2016guide} V.~Dumoulin and F.~Visin, ``A guide to convolution arithmetic for deep learning,'' \emph{arXiv preprint arXiv:1603.07285}, 2016.
\end{thebibliography}

\end{document}
